% =========================================================
% Calculus --- Limits
% =========================================================

\section{Limits of Functions}

% ---------------------------------------------------------
% Toolkit
% ---------------------------------------------------------

% ---------------------------------------------------------
\subsection*{The Definition}

\begin{propositionbox}{Definition (Limit of a Function)}
Let $f$ be defined near $a$, except possibly at $a$ itself.
The \emph{limit of $f(x)$ as $x$ approaches $a$} is $L$, written
\[
  \lim_{x \to a} f(x) = L,
\]
if for every $\varepsilon > 0$ there exists $\delta > 0$ such that
\[
  0 < |x - a| < \delta \implies |f(x) - L| < \varepsilon.
\]
\end{propositionbox}

\begin{remark*}[Intuition]
$f(x)$ can be made arbitrarily close to $L$ by taking $x$
sufficiently close to $a$ from either side, without $x$ reaching $a$.
The $\varepsilon$ names the required closeness of the output;
the $\delta$ responds with the required closeness of the input.
\end{remark*}

\begin{remark*}
The value $f(a)$ plays no role. The limit of $f(x) = \dfrac{x^2 - 1}{x - 1}$
as $x \to 1$ is $2$, even though the function is undefined at $x = 1$,
because $\dfrac{x^2-1}{x-1} = x + 1$ for all $x \neq 1$.
\end{remark*}


% ---------------------------------------------------------
\subsection*{Indeterminate Forms and Algebraic Techniques}

\begin{remark*}[Algebraic techniques]
Direct substitution fails when it produces $\tfrac{0}{0}$ or
$\tfrac{\infty}{\infty}$. The standard remedies are factoring,
rationalising, and cancellation.
\end{remark*}

\begin{tcolorbox}[colback=gray!6, colframe=gray!40, arc=2pt,
  left=6pt, right=6pt, top=4pt, bottom=4pt,
  title={\small\textbf{Example: Factoring}},
  fonttitle=\small\bfseries]
\[
  \lim_{x \to 1} \frac{x^2 - 1}{x - 1}
  = \lim_{x \to 1} \frac{(x-1)(x+1)}{x-1}
  = \lim_{x \to 1} (x + 1)
  = 2.
\]
\end{tcolorbox}

\begin{tcolorbox}[colback=gray!6, colframe=gray!40, arc=2pt,
  left=6pt, right=6pt, top=4pt, bottom=4pt,
  title={\small\textbf{Example: Rationalising}},
  fonttitle=\small\bfseries]
\[
  \lim_{x \to 0} \frac{\sqrt{x+1} - 1}{x}
  = \lim_{x \to 0} \frac{(\sqrt{x+1} - 1)(\sqrt{x+1} + 1)}{x(\sqrt{x+1}+1)}
  = \lim_{x \to 0} \frac{x}{x(\sqrt{x+1}+1)}
  = \frac{1}{2}.
\]
\end{tcolorbox}


% ---------------------------------------------------------
\subsection*{The Squeeze Theorem}

\begin{theorembox}{Theorem (Squeeze Theorem)}
Let $f$, $g$, and $h$ be functions on an open interval $I$
containing $c$ such that for all $x$ in $I$,
\[
  f(x) \leq g(x) \leq h(x).
\]
If
\[
  \lim_{x \to c} f(x) = L = \lim_{x \to c} h(x),
\]
then
\[
  \lim_{x \to c} g(x) = L.
\]
\end{theorembox}

\begin{remark*}[Intuition]
$g$ is trapped between $f$ and $h$.
If both outer functions converge to the same value $L$, the
middle function has nowhere else to go.
\end{remark*}

\begin{remark*}[Canonical application]
The standard use is $\lim_{x \to 0} \dfrac{\sin x}{x} = 1$,
established by bounding $\sin x / x$ between $\cos x$ and $1$
and applying the theorem as $x \to 0$.
\end{remark*}

% ---------------------------------------------------------
\subsection*{Limits of Functions Equal at All But One Point}

\begin{theorembox}{Theorem (Limits of Functions Equal at All But One Point)}
Let $g(x) = f(x)$ for all $x$ in an open interval containing $c$,
except possibly at $c$ itself, and let $\lim_{x \to c} g(x) = L$.
Then
\[
  \lim_{x \to c} f(x) = L.
\]
\end{theorembox}

\begin{remark*}
This theorem licenses the standard algebraic technique of
cancelling a common factor that vanishes at $c$ --- such as
the $(x-1)$ factor in $\dfrac{x^2-1}{x-1}$ --- and then
evaluating the simplified expression.
The original function and the simplified function agree
everywhere except at $c$, so they share the same limit there.
\end{remark*}

% ---------------------------------------------------------
\subsection*{One-Sided Limits}

\begin{propositionbox}{Definition (One-Sided Limits)}
\textbf{Left-hand limit.}
Let $f$ be defined on $(a, c)$ for some $a < c$.
The \emph{left-hand limit} of $f$ at $c$ is $L$, written
\[
  \lim_{x \to c^-} f(x) = L,
\]
if for any $\varepsilon > 0$ there exists $\delta > 0$ such that
for all $x \in (a, c)$, if $|x - c| < \delta$ then $|f(x) - L| < \varepsilon$.

\medskip
\textbf{Right-hand limit.}
Let $f$ be defined on $(c, b)$ for some $b > c$.
The \emph{right-hand limit} of $f$ at $c$ is $L$, written
\[
  \lim_{x \to c^+} f(x) = L,
\]
if for any $\varepsilon > 0$ there exists $\delta > 0$ such that
for all $x \in (c, b)$, if $|x - c| < \delta$ then $|f(x) - L| < \varepsilon$.
\end{propositionbox}

\begin{remark*}
The only difference from the two-sided definition is the
restriction of $x$ to one side of $c$.
One-sided limits are the natural tool for piecewise-defined
functions, where the formula changes at $c$.
\end{remark*}

% ---------------------------------------------------------
\subsection*{Limits and One-Sided Limits}

\begin{theorembox}{Theorem (Limits and One-Sided Limits)}
Let $f$ be defined on an open interval $I$ containing $c$,
except possibly at $c$. Then
\[
  \lim_{x \to c} f(x) = L
\]
if, and only if,
\[
  \lim_{x \to c^-} f(x) = L \quad \text{and} \quad
  \lim_{x \to c^+} f(x) = L.
\]
\end{theorembox}

\begin{remark*}
This theorem is the standard test for non-existence of a limit:
if the left- and right-hand limits exist but differ, the
two-sided limit does not exist.
The canonical example is $f(x) = |x|/x$ at $x = 0$, where
$\lim_{x \to 0^-} = -1$ and $\lim_{x \to 0^+} = 1$.
\end{remark*}
